\chapter{تصمیم‌گیری با شرط‌ها}%
\label{ch:conditionals}

تا اینجا برنامه‌های ما یک راستهٔ صاف بودند: خطِ اول، خطِ دوم، خطِ سوم، تمام.
اما «انتخاب»، دومین آجرِ سازندهٔ الگوریتم‌ها، به برنامه دوراهی می‌دهد: اگر
شرطی برقرار بود این کار را بکن، وگرنه آن کار را. در این فصل به برنامه‌هایمان
قدرتِ تصمیم‌گیری می‌دهیم.

\section{اگر: ساده‌ترین تصمیم}

\begin{salam}
تابع اصلی:
    سن : صحیح = 20
    اگر سن >= 18:
        چاپ "شما می‌توانید گواهینامه بگیرید"
    تمام
    چاپ "پایان برنامه"
تمام
\end{salam}

ساختارِ \code{اگر} چنین است:

\begin{center}
\code{اگر} \;\textit{شرط}\; \code{:} \quad\ldots\;بدنه\;\ldots\quad \code{تمام}
\end{center}

\emph{شرط} همان عبارت‌های منطقیِ فصلِ پیش است؛ چیزی که حاصلش \code{درست} یا
\code{نادرست} باشد. اگر شرط درست باشد، بدنه (خط‌های تورفتهٔ میانِ \code{اگر}
و \code{تمام}) اجرا می‌شود؛ اگر نادرست باشد، سلام از روی کلِ بدنه می‌پرد و
از خطِ بعد از \code{تمام} ادامه می‌دهد. در مثالِ بالا چون ۲۰ از ۱۸ بیشتر
است، هر دو پیام چاپ می‌شوند؛ اگر سن را ۱۵ بگذارید فقط «پایان برنامه» چاپ
می‌شود. خودتان امتحان کنید.

\begin{note}
خطِ «پایان برنامه» به هیچ شرطی وابسته نیست؛ چون بیرونِ بدنهٔ \code{اگر}
(بعد از \code{تمام}) قرار دارد. تورفتگی و \code{تمام} مرزِ بدنه را مشخص
می‌کنند؛ به این مرزها دقت کنید، چون جابه‌جا شدنِ یک خط به داخل یا خارجِ
بدنه، رفتارِ برنامه را عوض می‌کند.
\end{note}

\section{وگرنه: راهِ دوم}

بسیاری از تصمیم‌ها دو طرف دارند: اگر شرط برقرار بود این، \emph{وگرنه} آن.

\begin{salam}
تابع اصلی:
    عدد : صحیح = 37
    اگر عدد % 2 == 0:
        چاپ عدد, "زوج است"
    وگرنه:
        چاپ عدد, "فرد است"
    تمام
تمام
\end{salam}

\begin{progoutfa}
37 فرد است
\end{progoutfa}

این همان برنامهٔ زوج و فردِ فصلِ پیش است که حالا به‌جای \code{true} و
\code{false}، پیامِ آدمیزادی می‌دهد. دقت کنید که از دو بدنه، \emph{همیشه
دقیقاً یکی} اجرا می‌شود: یا بدنهٔ \code{اگر} یا بدنهٔ \code{وگرنه}؛ هرگز هر
دو، هرگز هیچ‌کدام.

\section{چند راهی: وگرنه با شرط}

گاهی بیش از دو حالت داریم. مثلاً تبدیلِ نمره به توصیف: ۱۷ به بالا «عالی»،
۱۲ به بالا «خوب»، ۱۰ به بالا «قبول» و پایین‌تر «مردود». برای این حالت،
بعد از \code{وگرنه} می‌توان شرطِ تازه‌ای آورد:

\begin{salam}
تابع اصلی:
    نمره : صحیح = 14
    اگر نمره >= 17:
        چاپ "عالی"
    وگرنه نمره >= 12:
        چاپ "خوب"
    وگرنه نمره >= 10:
        چاپ "قبول"
    وگرنه:
        چاپ "مردود"
    تمام
تمام
\end{salam}

\begin{progoutfa}
خوب
\end{progoutfa}

شرط‌ها \emph{به ترتیب از بالا} بررسی می‌شوند و اولین شرطِ درست، برنده است؛
بقیه حتی نگاه هم نمی‌شوند. نمرهٔ ۱۴ در شرطِ اول رد می‌شود (۱۴ کمتر از ۱۷
است)، در شرطِ دوم قبول می‌شود و «خوب» چاپ می‌شود؛ کاری به شرطِ سوم نمی‌رسد،
هرچند ۱۴ از ۱۰ هم بیشتر است. \code{وگرنه}ی آخر که شرطی ندارد، نقشِ «هیچ‌کدام
از بالایی‌ها» را بازی می‌کند.

\begin{warn}
ترتیبِ شرط‌ها مهم است. اگر همین برنامه را با شرطِ \code{نمره >= 10} در
\emph{بالا} بنویسید، نمرهٔ ۱۹ هم «قبول» می‌گیرد نه «عالی»؛ چون اولین شرطِ
درست همان‌جا پیدا می‌شود. قاعدهٔ سرانگشتی: در زنجیرهٔ شرط‌های عددی، از
سخت‌گیرترین شرط شروع کنید.
\end{warn}

\section{شرط درونِ شرط}

بدنهٔ \code{اگر} می‌تواند خودش یک \code{اگر}ِ دیگر داشته باشد. به این کار
\emph{تودرتو} می‌گویند. مثلِ این مثال که اول بررسی می‌کند عدد مثبت است یا
نه، و فقط برای مثبت‌ها زوج و فرد را می‌سنجد:

\begin{salam}
تابع اصلی:
    عدد : صحیح = -8
    اگر عدد > 0:
        اگر عدد % 2 == 0:
            چاپ "مثبت و زوج"
        وگرنه:
            چاپ "مثبت و فرد"
        تمام
    وگرنه:
        چاپ "مثبت نیست"
    تمام
تمام
\end{salam}

\begin{progoutfa}
مثبت نیست
\end{progoutfa}

هر \code{اگر} یک \code{تمام}ِ خودش را دارد؛ تورفتگیِ بیشترِ لایهٔ داخلی
کمک می‌کند جفت‌ها را گم نکنید. تودرتویی ابزارِ قدرتمندی است، اما اگر
لایه‌ها زیاد شوند خواندنِ برنامه سخت می‌شود؛ اغلب می‌توان با عملگرهای
منطقی (\code{\&\&} و \code{||}) یا زنجیرهٔ \code{وگرنه} همان منطق را
صاف‌تر نوشت. مثلاً «مثبت و زوج» را می‌شد یک‌راست نوشت:
\code{اگر عدد > 0 \&\& عدد \% 2 == 0:}.

\section{نگهبانِ تقسیم بر صفر}

قولِ فصلِ پیش را ادا کنیم. تقسیم بر صفر برنامه را متوقف می‌کرد؛ حالا
ابزارش را داریم که جلویش را بگیریم:

\begin{salam}
تابع اصلی:
    مجموع : صحیح = 0
    تعداد : صحیح = 0
    اگر تعداد != 0:
        چاپ "میانگین:", مجموع / تعداد
    وگرنه:
        چاپ "هنوز داده‌ای ثبت نشده است"
    تمام
تمام
\end{salam}

به چنین شرطی که پیش از یک عملیاتِ خطرناک می‌آید و حالتِ بد را جدا می‌کند،
\textbf{نگهبان} می‌گویند. برنامه‌نویس‌های خوب به حالت‌های مرزی (صفر، منفی،
خالی) بدبین‌اند و برایشان نگهبان می‌گذارند.

\section{اجرای دستی، این بار روی کد}

عادتِ اجرای دستیِ فصلِ \ref{ch:algorithms} را به کد هم بیاورید. برنامهٔ
نمره را با \code{نمره = 10} خط‌به‌خط دنبال کنید: شرطِ اول \code{10 >= 17}
نادرست؛ شرطِ دوم \code{10 >= 12} نادرست؛ شرطِ سوم \code{10 >= 10} درست؛
چاپ: «قبول». اگر پیش‌بینی‌تان با اجرای واقعی یکی بود، برنامه را
\emph{فهمیده‌اید}، نه فقط نوشته‌اید.

\section{تمرین‌ها}

\exercise برنامه‌ای بنویسید که متغیرِ \code{عدد} را بگیرد و بگوید «مثبت»،
«منفی» یا «صفر». هر سه حالت را آزمایش کنید.

\exercise الگوریتمِ بزرگ‌ترینِ سه عددِ فصلِ \ref{ch:algorithms} را حالا به
کد تبدیل کنید: سه متغیر بسازید و با \code{متغیر بزرگ‌ترین} و دو \code{اگر}،
بزرگ‌ترین را پیدا و چاپ کنید. با چند دستهٔ مختلف از مقدارها امتحانش کنید.

\exercise برنامه‌ای بنویسید که با گرفتنِ متغیرِ \code{سال}، کبیسه بودنش را
با پیامِ فارسی اعلام کند (عبارتِ منطقی‌اش را در تمرینِ فصلِ پیش ساختید).

\exercise ماشین‌حسابِ کوچک: سه متغیر بسازید: \code{الف} و \code{ب} (عدد) و
\code{عمل} (رشته‌ای شاملِ \code{"+"} یا \code{"-"} یا \code{"*"} یا
\code{"/"}). با زنجیرهٔ \code{اگر}/\code{وگرنه}، عملِ خواسته‌شده را انجام و
حاصل را چاپ کنید. برای تقسیم، نگهبانِ صفر یادتان نرود.

\exercise این برنامه اشکالی منطقی دارد؛ بدونِ اجرا پیدایش کنید:
\begin{salam}
تابع اصلی:
    دما : صحیح = 5
    اگر دما > 0:
        چاپ "هوا سرد است"
    وگرنه دما > 15:
        چاپ "هوا معتدل است"
    وگرنه:
        چاپ "یخبندان!"
    تمام
تمام
\end{salam}
با دماهای ۵، ۲۰ و ۵- خروجی چه می‌شود و چه باید می‌بود؟
